\completa{18}{Sergio}{5g: de Spark SQL a los seis libros Excel}
\senalar{Seguir las cuatro cajas; luego explicar cada opción visible en el fragmento y todas las filas de la tabla.}
\decir{El punto cinco g pide exportar las consultas cinco e y cinco f utilizando Spark. Por eso preguntamos: ¿cómo entregamos en Excel los resultados de las seis consultas de vistas temporales? Partimos de los resultados tabulares, para que puedan revisarse sus columnas y valores.}
\decir{El primer cuadro dice «spark punto sql; consultas cinco e y cinco f». Ese método ejecuta el texto SQL dentro de Spark. La flecha llega a «DataFrame de resultado; coalesce uno». Cada consulta devuelve un DataFrame. Aquí coalesce uno reduce sus particiones a una para la escritura. Esta llamada afecta la distribución del resultado; no reemplaza nulos como la función SQL de igual nombre. Es adecuada para estas tablas pequeñas, aunque concentrar una tabla grande en una sola partición puede limitar el rendimiento.}
\decir{La siguiente flecha conduce a «write punto format; proveedor spark excel». Write inicia la escritura; format selecciona el proveedor com punto crealytics punto spark punto excel, una biblioteca que permite generar archivos Excel desde Spark. La última flecha lleva a «seis libros punto xlsx; una hoja por consulta». Un libro es el archivo Excel completo y la hoja es la pestaña interna con la tabla. No son seis pestañas dentro de un único libro.}
\decir{En el fragmento, df es el DataFrame que recibe la función. La opción dataAddress construye el nombre de la hoja y la celda A uno como inicio de la tabla. Header true indica escribir los nombres de columna. InferSchema false es una opción del proveedor; en esta escritura el DataFrame ya tiene esquema y no necesita inferir tipos al leer. DateFormat indica el formato para fechas y timestampFormat añade hora, minuto y segundo para marcas de tiempo. Mode overwrite sustituye un archivo previo en el destino y save ruta escribe el resultado.}
\decir{La tabla de salida muestra cinco e uno, hoja cinco e uno variacion, con sesenta filas; cinco e dos, hoja cinco e dos perfil, con dieciocho; cinco e tres, hoja cinco e tres region, con nueve. Para cinco f uno, hoja cinco f uno distritos, hay diez; para cinco f dos, hoja cinco f dos nacionalidad, doce; y para cinco f tres, hoja cinco f tres meses, doce. Se guardan en la carpeta excel. Estos son nombres de hojas, no variables nuevas del análisis.}
\decir{El registro de ejecución documenta los seis libros escritos desde DataFrames. Hay una diferencia que debemos conocer: el export de cinco e dos conserva totales y promedio, pero no incluye la columna de porcentaje nacional de la consulta individual. Por eso, el cuarenta y dos coma cero dos por ciento se verifica en esa salida individual. Así cerramos la entrega de resultados y podemos volver a las preguntas que respondimos.}

\completa{19}{Sergio}{Conectar cada conclusión con su pregunta y evidencia}
\senalar{Sergio lee todas las filas y realiza el cierre. Terminar aquí si solo se realiza el recorrido principal.}
\decir{Para cerrar, volvemos a las preguntas. En tres b preguntamos cuántas personas contabilizadas hay por fila: el promedio es uno coma dos nueve cuatro uno y la mediana es uno. En tres c uno preguntamos qué peso tienen los menores: cincuenta y siete coma setenta y nueve por ciento nacional y setenta y dos coma veintiuno por ciento en Madre de Dios. En tres c dos preguntamos qué peso tienen los extranjeros: dos coma setenta y dos por ciento en 2024. Cada proporción tiene un denominador explícito.}
\decir{En tres c tres preguntamos qué meses superan la referencia uniforme: diciembre aporta nueve coma cuarenta y nueve por ciento, con índice uno coma catorce. En cinco e dos preguntamos qué perfil concentra más personas reportadas: adolescentes mujeres, con cincuenta y un mil veintiuna y cuarenta y dos coma cero dos por ciento nacional. Finalmente, en cuatro y cinco g preguntamos cómo conservar y entregar los resultados: la base y su respaldo mantienen ciento veintiún mil cuatrocientas veintisiete personas contabilizadas y el registro documenta seis libros Excel.}
\decir{El bloque inferior explica cómo se sostienen esas respuestas: cada ficha conecta el ítem, las columnas, las operaciones, el código original y la salida. Las veinticinco consultas pueden abrirse desde el mapa. Las conclusiones describen lo reportado en este dataset; las proporciones no estiman riesgo poblacional y los cambios temporales no demuestran causalidad. Con esto terminamos el recorrido principal. Muchas gracias; podemos abrir el detalle de cualquier consulta o realizar la demostración si disponemos de tiempo adicional.}
